\documentclass[11pt,a4paper]{article}
\usepackage[margin=25mm]{geometry}
\usepackage[T1]{fontenc}
\usepackage{lmodern,microtype}
\usepackage[hidelinks]{hyperref}
\setlength{\parindent}{0pt}
\setlength{\parskip}{9pt}
\pagestyle{empty}
\hypersetup{pdftitle={AIP ECC 2027 Proposal - Shivam Gupta},pdfauthor={Shivam Gupta}}
\begin{document}
\begin{center}
\textbf{7th AIP Early-Career Conference for Historians\\of the Physical Sciences}\\[5pt]
Proposal
\end{center}
\vspace{5pt}
\textbf{Name:} Shivam Gupta

\textbf{E-mail address:} \href{mailto:shivam1720406@gmail.com}{shivam1720406@gmail.com}

\textbf{Institutional affiliation:} Independent Scholar

\textbf{Presentation title:} After Krypton: Realization, Uncertainty, and the Uneven Remaking of the Metre, 1983-2000

\textbf{Abstract}

In 1983 the krypton-based definition of the metre was abolished, but krypton remained an authorized realization. How did international metrology change a unit without synchronizing the retirement of its instruments? I compare CGPM decisions, the 1963 secondary-wavelength recommendation, the CIPM radiation lists of 1983, 1992 and 1997, and BIPM and Canadian laboratory reports. A reproducible reconstruction distinguishes 31 assigned regimes and twelve inherited lamp lines, whose wavelengths recur unchanged between 1963 and 1997. Preserving uncertainty conventions reduces the apparent 1983-1992 improvement of the 633 nm recommendation from fortyfold to 13.6-fold; tighter operating conditions further qualify what that ratio means. Canadian records show source replacement proceeding around calibration services, with image processing among its motivations. I argue that redefinition redistributed authority through revisable, conditional recommendations rather than imposing a common timetable of replacement. Linking numerical prescriptions to repair and staged modernization reveals continuity as an institutional and material achievement, and offers historians a method for tracing technical change beyond successive definitions.

\textbf{Short biography}

Shivam Gupta is an early-career independent scholar and software engineer with a Bachelor of Technology in Computer Science and Design from the Indraprastha Institute of Information Technology Delhi. His background includes undergraduate research and professional work in software engineering and product development. His current research brings reproducible computational methods to the history of measurement, with particular attention to standards, instrumentation, maintenance and the interpretation of technical records. He is pursuing this project independently and does not hold a PhD.
\end{document}
